\documentclass[10pt]{article}

\usepackage[utf8]{inputenc}

\usepackage[T1]{fontenc}

\usepackage{amsmath}

\usepackage{amsfonts}

\usepackage{amssymb}

\usepackage[version=4]{mhchem}

\usepackage{stmaryrd}

\usepackage{graphicx}

\usepackage[export]{adjustbox}

\graphicspath{ {./images/} }



\begin{document}

Аналогично определяется Т. м. для абелева многообразия. Пусть $A$ - абелево многообразие, определенное над $k$ и $A_{p^{n}}-$ группа точек порядка $p^{n}$ в $A(\bar{k})$. Тогда $T_{p}(A)$ определяется как $\lim _{p n}$. М о д у л е м Т е й-



\includegraphics[max width=\textwidth, center]{2023_07_09_c55401158e3ae4ad009eg-1(2)}

этой кривой.



Конструкция модуля $T_{p}(X)$ обобщается на случай числовых полей. Пусть $k$ - поле алгебраич. чисел и $k_{\infty}$ - нек-рое $Z_{p}$-расширение поля $k$ (расширение $\mathrm{c}$ группой Галуа изоморфной $Z_{p}$ ). Для промежуточного поля $k_{n}$ степени $p^{n}$ над $k$ пусть $\mathrm{Cl}\left(k_{n}\right)_{p}$ есть $p$-компонента группы классов идеалов поля $k_{n}$. Тогда $T_{p}\left(k_{\infty}\right)=$ $=\underline{\longleftarrow} \mathrm{Cl}\left(k_{n}\right)_{p}$, где предел берется относительно норменных отображений $\mathrm{Cl}\left(k_{m}\right)_{p} \rightarrow \mathrm{Cl}\left(k_{n}\right)_{p}$ для $m>n$. Модуль $T_{p_{0}}\left(k_{\infty}\right)$ характеризуется своими и в в а p и а нт а м и И в а с а в ы $\lambda, \mu, v$, к-рые определяются из условия



\[

\left|\mathrm{Cl}\left(k_{n}\right)_{p}\right|=p^{e_{n}}

\]



где $e_{n}=\lambda n+\mu p^{n}+v$ для всех достаточно больших $n$. Имеется предположение, что для круговых $Z_{p}$-расширений инвариант $\mu$ равен 0 . Это доказано для абелевых полей [4]. Известны примеры некруговых $Z_{p}$-расшширений с $\mu>0$ (см. [3]). Даже в случае, когда $\mu=0$, $T_{p\left(k_{\infty}\right) \text { не обязан быть свободным }} Z_{p \text {-модулем. }}$



Лит.: [1] Т э й т Д ж., «Математика», 1969, т. 13, №2, с. $3-25 ;$ [2] ші а ф а р е в и ч И. Р. б-функция, М., 1969; and commutative algebra, Tokyo, 1973, p. $1-11 ;$ [4] F e r r er o B., W a s h i ng t o n L. C., "Ann. Math.", 1979, v. 109,

p. 377-95.



ТЕJЕГРАФНОЕ УРАВНЕНИЕ - диффференциальное уравнение с тастными производными



\[

\frac{\partial^{2} u}{\partial t^{2}}-c^{2} \frac{\partial^{2} u}{\partial s^{2}}+(\alpha+\beta) \frac{\partial u}{\partial t}+\alpha \beta u=0 .

\]



Этому уравнению удовлетворяет напряжение тока в проводе, рассматриваемое как функция времени $t$ и расстояния $s$ от нек-рой фиксированной точки провода. Здесь $c$ - скорость света, $\alpha$ - емкостный, $\beta$ - индуктивный коэфффициенты .



Преобразованием



$e^{1 / 2(\alpha+\beta) t} u(s, t)=v(x, y), x=s+c t, y=s-c t$,



уравнение (1) приводится к виду



\[

v_{x y}+\lambda v=0, \quad \lambda=\left(\frac{\alpha-\beta}{4 c}\right)^{2} .

\]



Это уравнение принадлежит к классу гиперболич. уравнений 2-го порядка



\[

v_{x y}+a v_{x}+b v_{y}+c v=f,

\]



в теории к-рых важную роль играет функция Римана $R(x, y ; \xi, \eta)$. Для уравнения (2) эта функция выписывается в явном виде:



\[

R(x, y ; \xi, \eta)=J_{0}(\sqrt{4 \lambda(x-\xi)(y-\eta)}),

\]



где $J_{0}(z)-$ функция Бесселя.



Лит.: [1] К у р а н т Р., Уравнения с частными производными, пер. с англ., М., 1964.



T E JI E C H GI И УГОЛ-тасть пространства, ограниченная одной полостью нек-рой конич. поверхности (рис.), направляющая к-рой гомеоморфна окружности. Частными случаями Т. У. яв-



\includegraphics[max width=\textwidth, center]{2023_07_09_c55401158e3ae4ad009eg-1(1)}

A. II. Солдатов ляются многогранные углы. За меру Т. у. принимают отношение площади той части сферы с центром в вер- шине конич. поверхности, к-рая вырезается этим Т. у., к квадрату радиуса сферы. Напр., Т. у., заключающий $1 / 8$ часть пространства (октант), измеряется числом $4 \pi R^{2} / 8 R^{2}=\pi / 2$. Единицей измерения Т. $\mathbf{~}$. является стерадиан.



БCЭ-3.



TЕЛО - кольцо, в к-ром уравнения $a x=b$ и $y a=b$, где $a \neq 0$, однозначно разрешимы. В случае ассоциативного кольца достаточно потребовать существования единицы 1 и однозначной разрешимости уравнений $a x=1$ и $y a=1$ для любого $a \neq 0$. Коммутативное ассоциативное Т. является полем. Пример некоммутативного ассоциативного Т. - т е л о к в а т е рн и о н о в, определяемое как множество матриц вида $\left\|\begin{array}{r}a \bar{b} \\ -b \bar{a}\end{array}\right\|$ над полем комплексных чисел с обычными операциями (ср. Кватернион). Примером неассоциативного Т. является Кәли - Диксона аләебра, состоящая из всех матриц того же вида над телом кватернионов. Это Т. альтернативно (см. Альтернативнье кольца и алгебрь). Всякое Т. является алгеброй с делением или над полем рациональных чисел, или над полем вычетов. Тело кватернионов является 4-мерной алгеброй над полем действительных чисел, а алгебра Кәли - Диксона 8-мерной. Размерность любой алгебры с делением над полем действительных чисел равна 1, 2, 4 или 8 (cм. [1], а также Топологическое кольцо). Поля действительных и комплексных чисел и тело кватернионов и только они являются связными локально компактными ассоциативными Т. (см. [5]). Всякая конечномерная алгебра без делителей нуля есть Т. Всякое конечное ассоциативное T. коммутативно (см. [6], [8]). Ассоциативное Т. характеризуется тем, что все ненулевые модули над ним свободны. Всякое неассоциативное альтернативное Т. конечномерно [3]. Сходный результат верен для тел Мальцева [7] (см. Мальцева алгебра) и йордановых Т. [4] (см. Порданова алгебра). В отличие от коммутативного случая, не всякое ассоциативное кольцо без делителей вуля вложимо в Т. (см. Вложение кольца).



с. Лит.: [1] А д а м с Д. Ф., «Математика", 1961, т. 5, № 4, с нем., 2 изд., М., 1979; [3] Ж е в л а к о в К. А. и пр., Кольца, бллизкие к ассоциативным, М., 1978; [4] З е л ь м ано в Е. И., "Алгебра и логика", 1979, т. 18, Nㅜ 3, с. 286-310; [5] П о н т р я г и н Л. С., Непрерывные групшы, 3 изд., М., 1973; [6] С. к о р н я к о в Л. А., Элементы общей алтебры, М., 1983 ; [7] Ф и Л и П П о в В. Т., “Алгебра и логика», 1976, т. 15,

$\mathcal{N}_{2} 2$, с. $235-42 ;$ [8] Х е р с т е й И., Некоммутативные кольца, пер. с англ., М., $1972 . \quad$ Л. А. Скорняков.



ТEHЗОР на век т о р ном п р о с т р а н т в $V$ на д по л е м $k$ - элемент $t$ векторного пространства



\[

T^{p, q}(V)=(\stackrel{p}{\otimes} V) \otimes\left(\stackrel{q}{\otimes} V^{*}\right),

\]



где $V^{*}=\operatorname{Hom}(V, k)$ - пространство, сопряженное с $V$. Говорят, что тензор $t$ является $p$ раз контравариантным и $q$ раз ковариантным или что $t$ имеет тип $(p, q)$. Число $p$ наз. К о н т р а в а р и а н т о й в а л е н т но-



\includegraphics[max width=\textwidth, center]{2023_07_09_c55401158e3ae4ad009eg-1}

с т ь ю, а число $p+q-06$ е й в а ленте о с тью тензора $t$. Пространство $T^{0,0}(V)$ отождествляется с $k$. Тензоры типа $(p, 0)$ наз. к о н т $\mathbf{p}$ а в а $\mathbf{p}$ и а н т ным и, типа $(0, q)$ - к о в а р и а в т ны м и, а остальные - с м е ї а н н ы ми.



П р и м е р Ы Т. 1) Вектор пространства $V$ (Т. типа $(1,0))$.



\begin{enumerate}

  \setcounter{enumi}{1}

  \item Ковектор пространства $V$ (Т. типа $(0,1))$.



  \item Каждый ковариантный $\mathrm{T}$.



\end{enumerate}



\[

t=\sum_{i=1}^{s} h_{i 1} \otimes \ldots \otimes h_{i q},

\]



где $h_{i j} \in V^{*}$, определяет $q$-линейную форму $\hat{t}$ на $V$ по формуле



\[

\hat{t}\left(x_{1}, \ldots, x_{q}\right)=\sum_{i=1}^{s} h_{i 1}\left(x_{1}\right) \ldots h_{i q}\left(x_{q}\right)

\]





\end{document}